% Page budget: 1.0 page. Demand response and $/MWh; Philly facts (failed jobs die early, many killed, daily power shape); F2.
% S10: every number comes from numbers.tex (macros \Mot..., written by paper/figures/motivation.py from results/motivation/stats.csv).
\section{Background and Motivation}
\label{sec:background}

\subsection{Demand response and dollars per MWh}
Electric grids must balance supply and demand at every moment, and operators ask large consumers to reduce their load for a few hours when demand is high. This practice is called demand response~\cite{liu13}. A data center can take part because part of its load is flexible in time: batch work can wait, and equipment can run slower. Earlier work studies how data centers can avoid a coincident peak by shifting workload and using local generation~\cite{liu13}, and compares the response potential of GPU and CPU data centers~\cite{zhou24}. Carbon-aware computing uses the same flexibility to follow a cleaner grid~\cite{radovanovic23,wiesner21}.

A request has a start, a length and a size: reduce power by a fraction for a given number of hours (Section~\ref{sec:problem}). A cluster should compare not only whether it can comply, but at what price. We therefore express every action as the dollars spent for each megawatt-hour of load actually removed. Delaying a job costs the value of the delayed GPU-hours, capping power slows work, and renting cloud GPUs costs the cloud bill. We use dollars per MWh because it puts the three actions on one energy-based scale and, unlike a count of delayed jobs, it exposes an action that sheds little energy for much money.

\subsection{What the Philly trace shows}
We use the Microsoft Philly trace~\cite{jeon19}, which records deep-learning jobs on a shared multi-tenant GPU cluster. The test window contains \MotTestJobs{} jobs and \MotTestGpuHours{} GPU-hours of work, which our simulator runs on a pool of \MotTestPoolGpus{} GPUs. Figure~\ref{fig:motivation} shows four facts that shape the problem.

\begin{figure*}[t]
\centering
\includegraphics[width=\textwidth]{figures/out/f2_motivation.pdf}
\caption{The Philly test window. (a)~Share of jobs and of GPU-hours by final status. (b)~Cumulative distribution of wall time by status (logarithmic axis; the grey line marks one hour). (c)~Baseline facility power over one week; the shaded bands are 17:00--20:00 Pacific time. (d)~Histogram of per-job mean GPU utilisation; the dashed line is the minute-weighted mean (\MotUtilMeanMinuteweightedPct\%).}
\label{fig:motivation}
\end{figure*}

\emph{Failed and killed jobs are a large share of the work.} \MotJobsharePassPct\% of jobs end as Pass and use \MotGpuhsharePassPct\% of the GPU-hours, whereas Failed jobs are \MotJobshareFailedPct\% of jobs and \MotGpuhshareFailedPct\% of GPU-hours, and Killed jobs are \MotJobshareKilledPct\% of jobs and \MotGpuhshareKilledPct\% of GPU-hours (Figure~\ref{fig:motivation}a). Together, Failed and Killed jobs consume \MotGpuhshareFailkillPct\% of the GPU-hours. The original analysis of the trace also reports failures during training~\cite{jeon19}.

\emph{Runtimes spread over orders of magnitude and differ by status.} The median wall time is \MotWallMedianPassMin{} min for Pass, \MotWallMedianFailedMin{} min for Failed and \MotWallMedianKilledMin{} min for Killed jobs, while the 90th percentile is \MotWallPNineZeroPassMin{}, \MotWallPNineZeroFailedMin{} and \MotWallPNineZeroKilledMin{} minutes, respectively (Figure~\ref{fig:motivation}b). \MotWithinOnehFailkillPct\% of Failed and Killed jobs end within one hour of their first start, yet the rest run longer.

\emph{Cluster power is flat, not peaked.} In the baseline simulation, facility power averages \MotPowerMeanKw{} kW with a highest value of \MotPowerPeakKw{} kW. The mean over the 17:00--20:00 Pacific-time windows is \MotPowerEveningMeanKw{} kW, a ratio of \MotPowerEveningOverOverall{} to the overall mean (Figure~\ref{fig:motivation}c). The evening is therefore a peak of the grid, not of the cluster: the cluster is asked to reduce an ordinary load. In our power model, the pool draws at least \MotPowerIdleFloorKw{} kW even when every GPU is idle, so no action can cut deeper than this floor.

\emph{GPUs are lightly used on average.} The mean per-job GPU utilisation is \MotUtilMeanJobweightedPct\% over jobs and \MotUtilMeanMinuteweightedPct\% when weighted by measured minutes, with a median of \MotUtilMedianPct\% (Figure~\ref{fig:motivation}d). Power is therefore not proportional to the number of busy GPUs, which is why the power model of Section~\ref{sec:problem} uses each job's utilisation.

\subsection{Why runtime and failure predictions matter}
When the grid calls, the cluster must decide which running jobs to touch. Suspending a job delays useful work only if the job would have produced useful output; delaying a job that will fail anyway costs little. Bursting a job pays the cloud for its whole remaining runtime, so it is attractive only for a job that will end soon and succeed. Both decisions need an estimate of the remaining runtime and of the chance of failure at the moment of the event.

Runtime prediction for scheduling has a long history. User-supplied estimates are poor, and predictions generated by the system from a user's recent jobs improve backfilling~\cite{tsafrir07}. For GPU clusters, learned predictors use features of the job and of the user's history~\cite{hu21,hu23}. Work on hybrid clouds decides when to send jobs to rented capacity to trade cost against completion time~\cite{assuncao09,luo24}; it does not treat offloading as a priced grid action.

The jobs running at an event are not typical jobs. At the \MotEventStarts{} event starts of our experiments, a mean of \MotRunningJobsAtStartsMean{} jobs are running, and \MotRunningBeyondGridJobsPct\% of them (\MotRunningBeyondGridGpusPct\% of the running GPUs) have already run for more than 12 hours, which is beyond the last elapsed time at which our predictors give a value. Predictions for long-running jobs therefore need care, which we return to in Section~\ref{sec:design}.
